\section{系统约束与稳定性保证}
\subsection{资源与算力限界}
策略梯度训练常面对有限的算力与 budget：长 horizon 轨迹意味着巨大的 GPU 时间，而 entropy regularization 需要额外 compute。我们通常通过 accumulate gradients、梯度剪裁与 mixed precision 来平衡训练负载，并在 pipeline 中加入时序 dependency 检查防止 drift。

\begin{table}[H]
\centering
\small
\begin{tabular}{p{0.34\textwidth}p{0.38\textwidth}p{0.28\textwidth}}
\toprule
约束类型 & 具体表现 & 工程应对 \\
\midrule
延迟/预算 & 轨迹 horizon 长导致 compute burst & gradient accumulation + distributed rollout \\
稳定性 & entropy collapse 造成 policy collapse & entropy bonus + KL guard \\
数据可用性 & limited rollout data per day & replay buffer + prioritized sampling \\
\bottomrule
\end{tabular}
\caption{系统约束下策略梯度的资源格局}
\end{table}

\begin{importantbox}{算力与 logging 的联动}
\begin{itemize}
    \item 每个 rollout 阶段都要写入 seed, config, entropy curve 以便复现；
    \item 长 horizon 轨迹在记录时需 chunking，避免 log 文件爆炸；
    \item Closed-loop pipeline 里把 gradient step count 与 runtime metric 绑定，便于 capacity planning。
\end{itemize}
\end{importantbox}

\subsection{制度与公平}
除了硬件，模型部署还受到合规、公平与安全的制度约束。例如在 multi-domain rollout 中必须检测 policy 对不同群体的 reward 分布，如果差距过大需要 trigger guardrail。此处常用 fairness dashboard 记录 disparity, reject rate 与 manual override logs。

\begin{warningbox}{忽略制度要求的风险}
如果只关注 reward 而忽略 fairness 指标，policy 可能在 release 后因差异化表现被 regulator 叫停，延缓整个产品进度。
\end{warningbox}

\begin{knowledgebox}{制度化保障的做法}
\begin{itemize}
    \item 设立 `fairness gate`：对敏感 subgroup 的 reward/distribution 做提前评估；
    \item 对 `override event` 进行 structured logging，关联 policy version 与 prompt template；
    \item 用 audit-ready benchmark（如 fairness-specific tasks）确保多维度对齐。
\end{itemize}
\end{knowledgebox}

\subsection{本章小结}
系统约束既来自算力 budget，也来自公平与 compliance 要求。通过细化 logging、entropy guard 以及 fairness gate，可以把策略梯度训练拉回到可管理的工程轨道。

